% 子集（综述）
% license CCBYSA3
% type Wiki

本文根据 CC-BY-SA 协议转载翻译自维基百科\href{https://en.wikipedia.org/wiki/Subset}{相关文章}

\begin{figure}[ht]
\centering
\includegraphics[width=6cm]{./figures/d095c5b69171d4da.png}
\caption{欧拉图示意：A 是 B 的子集（记作 $A \subseteq B$），反之，B 是 A 的超集（记作 $B \supseteq A$）。} \label{fig_ZiJ_1}
\end{figure}
在数学中，如果集合A的所有元素也是集合B的元素，则称A是B的子集；此时$B$是$A$的超集。集合$A$与$B$有可能相等；如果它们不相等，则A是B的真子集。一个集合是另一个集合的子集这一关系称为包含关系（有时也称为包涵关系）。$A$是$B$的子集也可以表述为$B$包含$A$或$A$被包含于$B$。一个k-子集是指具有k个元素的子集。

在量化表示时，$A \subseteq B$可写作$\forall x \,(x \in A \Rightarrow x \in B)$。

可以通过一种称为元素论证法的证明技巧来证明命题$A \subseteq B$。
\begin{enumerate}
\item 设给定集合$A$和$B$。要证明$A \subseteq B$。
\item 假设a是集合A中一个特定但任意选择的元素，并证明a也是集合B的元素。
\end{enumerate}
这种方法的有效性可以看作是全称概化的结果：该方法表明，对于任意选择的元素c，都有
$(c \in A) \Rightarrow (c \in B)$。由此，全称概化推出$\forall x \,(x \in A \Rightarrow x \in B)$，这等价于$A \subseteq B$如上所述。
\subsection{定义}
如果$A$和$B$是集合，且$A$的每一个元素也是$B$的元素，那么：
\begin{itemize}
\item $A$是$B$的子集，记作$A \subseteq B$或等价地，
\item $B$是$A$的超集，记作$B \supseteq A$
\end{itemize}
如果A是B的子集，但A不等于B（即至少存在一个B中的元素不在A中），那么：
\begin{itemize}
\item $A$是$B$的真（或严格）子集，记作$A \subsetneq B$或等价地，
\item $B$是$A$的真（或严格）超集，记作$B \supsetneq A$。
\end{itemize}
空集，写作$\{\}$或$\varnothing$，没有任何元素，因此平凡地是任意集合$X$的子集。
\subsection{基本性质}
\begin{itemize}
\item 自反性：对任意集合$A$都有$A \subseteq A$
\item 传递性：若$A \subseteq B$且$B \subseteq C$则$A \subseteq C$
\item 反对称性：若$A \subseteq B$且$B \subseteq A$则$A = B$
\end{itemize}
\subsubsection{真子集}
反自反性：对任意集合$A$,$A \subsetneq A$为假。传递性：若$A \subsetneq B$且$B \subsetneq C$,则$A \subsetneq C$

**反对称性（Asymmetry）**：若$A \subsetneq B

则

$$
B \subsetneq A
$$

为假。
